% \iffalse meta-comment
%
% Copyright (C) 2026 Evan Berkowitz
%
% This work may be distributed and/or modified under the conditions of the
% LaTeX Project Public License, either version 1.3c of this license or (at
% your option) any later version.  The latest version of this license is in
%   https://www.latex-project.org/lppl.txt
%
% This work has the LPPL maintenance status `maintained'.
% The Current Maintainer of this work is Evan Berkowitz.
% This work consists of the files mermin.dtx and mermin.ins and the derived
% file mermin.sty.
%
% \fi
%
% \iffalse
%<*driver>
\ProvidesFile{mermin.dtx}
%</driver>
%<package>\NeedsTeXFormat{LaTeX2e}[2020/10/01]
%<package>\ProvidesPackage{mermin}
%<*package>
    [2026/11/16 v0.1 Mermin-style equation tools]
%</package>
%
%<*driver>
\documentclass{ltxdoc}
\usepackage{amsmath}
\usepackage{mermin}
\usepackage[hidelinks]{hyperref}
\EnableCrossrefs
\CodelineIndex
\RecordChanges
\begin{document}
  \DocInput{mermin.dtx}
  \PrintChanges
  \PrintIndex
\end{document}
%</driver>
% \fi
%
% \CheckSum{0}
%
% \changes{v0.1}{2026/11/16}{Initial version, providing \cs{alsolabel}.}
%
% \GetFileInfo{mermin.sty}
%
% \DoNotIndex{\begingroup,\endgroup,\def,\gdef,\let,\global,\else,\fi}
% \DoNotIndex{\ifx,\ifdefined,\expandafter,\noexpand,\unexpanded,\xdef}
% \DoNotIndex{\newcommand,\newif,\label,\do,\the,\toks@,\@xp,\@empty}
% \DoNotIndex{\AddToHook,\NeedsTeXFormat,\ProvidesPackage,\RequirePackage}
% \DoNotIndex{\@ifnextchar,\endinput}
%
% \title{The \textsf{mermin} package\thanks{This document corresponds to
%   \textsf{mermin}~\fileversion, dated \filedate.}}
% \author{Evan Berkowitz}
% \date{\filedate}
%
% \maketitle
%
% \begin{abstract}
% The \textsf{mermin} package helps physicists write equations the way
% N.~David Mermin asks in ``What's Wrong with These Equations?''\
% \cite{mermin}.
% Its first tool, \cs{alsolabel}, gives one numbered equation, or one line
% of an \textsf{align}, as many labels as the prose needs, so that every
% citation of that equation can carry a key that matches the noun phrase
% used to cite it.
% \end{abstract}
%
% \section{Introduction}
%
% In his 1989 column~\cite{mermin} Mermin proposed three rules for
% equations in physics writing.
% Fisher's rule asks that every displayed equation be numbered, whether or
% not the text refers to it yet, because readers, referees and later papers
% will.
% The Good Samaritan rule asks that an equation be cited by a phrase as well
% as a number, so that the reader need not page back to learn what
% ``(11)'' says: one writes ``the anomaly condition~(11)'', never a bare
% ``(11)'' and never ``Eq.~(11) shows''.
% An equation number locates; the noun phrase says what is there.
% The third rule, that mathematics is prose, asks that a display be
% punctuated as part of its sentence, so that the paragraph reads straight
% through the equation.
%
% The Good Samaritan rule has a cost that \LaTeX\ makes worse than it need
% be.
% The right phrase for an equation depends on why it is being cited.
% The same line may be ``the anomaly condition'' in one paragraph and
% ``the Pythagorean condition'' in another, yet \textsf{amsmath} allows a
% numbered line exactly one \cs{label}.
% The key then names only one of the roles, and the source no longer says
% what the author meant at each citation.
% The \cs{alsolabel} command removes this restriction without changing
% \cs{label}.
%
% \section{Usage}
%
% \DescribeMacro{\alsolabel}
% Load the package with |\usepackage{mermin}|.
% Then |\alsolabel{|\meta{key}|}| attaches \meta{key} to the current
% equation, or to the current line of a multi-line display, in addition
% to its \cs{label}.
% It may be used any number of times on a line, before or after the
% \cs{label}, or on a line that has no \cs{label} at all.
% Every key refers to the same number and, with \textsf{hyperref}, to the
% same anchor.
% Consider two-flavor chiral gauge theories whose gauge anomaly cancels
% when the left- and right-handed charges obey
% \begin{equation}
%   L_1^2+L_2^2=R_1^2+R_2^2.
%   \label{eq:anomaly}
%   \alsolabel{eq:pythagorean}
% \end{equation}
% The source of this display reads
% \begin{verbatim}
% \begin{equation}
%   L_1^2+L_2^2=R_1^2+R_2^2.
%   \label{eq:anomaly}
%   \alsolabel{eq:pythagorean}
% \end{equation}
%
% The 3450 model, with charges $(3,4\,\vert\,5,0)$, satisfies the anomaly
% condition~\eqref{eq:anomaly}.
% Read as a statement about integers, the same line is the Pythagorean
% condition~\eqref{eq:pythagorean}, and $3^2+4^2=5^2$ is its smallest
% nontrivial solution.
% \end{verbatim}
% which gives each citation a key that matches its phrase and typesets as
% \begin{quote}
% The 3450 model, with charges $(3,4\,\vert\,5,0)$, satisfies the anomaly
% condition~\eqref{eq:anomaly}.
% Read as a statement about integers, the same line is the Pythagorean
% condition~\eqref{eq:pythagorean}, and $3^2+4^2=5^2$ is its smallest
% nontrivial solution.
% \end{quote}
%
% Lines of an \textsf{align} are labeled the same way, one set of keys per
% line, and a line may be labeled by \cs{alsolabel} alone:
% \begin{verbatim}
% \begin{align}
%   a &= b  \label{eq:first} \alsolabel{eq:first-again} \\
%   c &= d  \alsolabel{eq:second}
% \end{align}
% \end{verbatim}
% Outside a display \cs{alsolabel} is simply \cs{label}, so it can also
% give a section or a figure a second name.
%
% Because \cs{label} is left alone, two \cs{label}s on one line remain the
% error \textsf{amsmath} has always reported; only \cs{alsolabel} adds
% keys.
% Citing a key with \cs{ref}, \cs{eqref}, \cs{pageref} or the
% \textsf{cleveref} command \cs{cref} works exactly as for a key given to
% \cs{label}.
%
% \section{Compatibility}
%
% The package works in the \textsf{amsmath} environments \textsf{equation},
% \textsf{align}, \textsf{gather}, \textsf{multline}, \textsf{split} (with
% the keys given in the enclosing \textsf{equation}) and
% \textsf{subequations}, and in \textsf{empheq}.
% It works with \textsf{hyperref} and \textsf{cleveref} loaded either
% before or after it.
% It requires a \LaTeX\ format from 2020-10-01 or later, for the
% |begindocument/end| hook.
% The \textsf{mathtools} option |showonlyrefs| has not been tested.
%
% The repository that holds this package also holds a test suite.
% Running |make test| typesets every environment above, with the keys
% given in every order, in four package configurations, and checks that
% each key in the resulting \texttt{.aux} file carries its line's number
% and anchor.
%
% \begin{thebibliography}{1}
% \bibitem{mermin}
%   N.~D. Mermin, ``What's wrong with these equations?'',
%   \emph{Physics Today} \textbf{42}(10), 9 (1989),
%   \href{https://doi.org/10.1063/1.2811173}{doi:10.1063/1.2811173}.
% \end{thebibliography}
%
% \StopEventually{}
%
% \section{Implementation}
%
% \textsf{amsmath} does not write a display label when it meets \cs{label}.
% It stores the key in |\df@label| and writes it with
% |\ltx@label{|\meta{contents of \texttt{\char`\\df@label}}|}| once the
% line's tag, and so |\@currentlabel| and the \textsf{hyperref} anchor, are
% known.
% The package keeps the extra keys of a line in a queue and writes them
% immediately after that primary label, so they receive the same number and
% anchor.
%
% Three complications shape the code.
% First, \textsf{cleveref} stores keys in a different form, wrapping the key
% in braces or preceding it with an optional argument, so every queued key
% is stored in exactly the form |\label@in@display| would have stored it.
% Second, \textsf{amsmath} and \textsf{empheq} typeset some displays twice,
% once to measure them, and may discard a primary label without writing it.
% The queue is therefore tagged with the primary label it belongs to and
% written only when that label is, and a key is never queued twice.
% Third, \cs{alsolabel} may come before \cs{label} or without it.
% It then occupies |\df@label| itself, and a later \cs{label} reclaims the
% slot and moves the earlier key into the queue.
%
%    \begin{macrocode}
%<*package>
\RequirePackage{amsmath}
%    \end{macrocode}
%
% \begin{macro}{\mermin@queue}
% \begin{macro}{\mermin@tag}
% \begin{macro}{\mermin@promoted}
% The queue holds extra keys as |\do{|\meta{stored form}|}|.
% The tag is the |\df@label| the queue belongs to.
% |\mermin@promoted| is nonempty when \cs{alsolabel} occupies |\df@label|
% because no \cs{label} has come yet.
%    \begin{macrocode}
\gdef\mermin@queue{}
\gdef\mermin@tag{}
\gdef\mermin@promoted{}
\newif\ifmermin@also
%    \end{macrocode}
% \end{macro}
% \end{macro}
% \end{macro}
%
% \begin{macro}{\alsolabel}
% The user command calls whatever \cs{label} means at that point, with a
% flag raised.
% The measuring passes of \textsf{amsmath} gobble or defer \cs{label}, and
% \cs{alsolabel} thereby inherits the same treatment.
% In a display \cs{label} is |\label@in@display|, which the package wraps
% below to read the flag.
% \textsf{cleveref}'s measuring pass for \textsf{multline} skips
% |\label@in@display| and calls \textsf{amsmath}'s original, saved as
% |\cref@old@label@in@display|, directly with the stored form of the key;
% that hook is redirected as well.
%    \begin{macrocode}
\newcommand\alsolabel[1]{%
  \begingroup
    \mermin@alsotrue
    \ifdefined\mermin@cref@old
      \let\cref@old@label@in@display\mermin@also@stored
    \fi
    \label{#1}%
  \endgroup}
%    \end{macrocode}
% \end{macro}
%
% \begin{macro}{\mermin@call@orig}
% This calls the |\label@in@display| that was in force before the package
% wrapped it, with \textsf{cleveref}'s inner hook restored.
% The argument is the complete argument list, braces included.
%    \begin{macrocode}
\def\mermin@call@orig#1{%
  \begingroup
    \ifdefined\mermin@cref@old
      \let\cref@old@label@in@display\mermin@cref@old
    \fi
    \mermin@orig@label@in@display#1%
  \endgroup}
%    \end{macrocode}
% \end{macro}
%
% \begin{macro}{\mermin@canonical}
% This sets |\mermin@form| to the form in which |\label@in@display| would
% store the key |#1|, by letting it store the key into an emptied
% |\df@label| and restoring the old contents afterwards.
%    \begin{macrocode}
\def\mermin@canonical#1{%
  \global\let\mermin@saved\df@label
  \global\let\df@label\@empty
  \mermin@call@orig{{#1}}%
  \global\let\mermin@form\df@label
  \global\let\df@label\mermin@saved
}
%    \end{macrocode}
% \end{macro}
%
% \begin{macro}{\mermin@enqueue}
% This queues |\mermin@form| for the current |\df@label|.
% A queue tagged with a different label is stale and is discarded first.
% A form already in the queue is skipped, because a display typeset twice
% runs its \cs{alsolabel}s twice.
%    \begin{macrocode}
\def\mermin@enqueue{%
  \ifx\mermin@tag\df@label\else
    \gdef\mermin@queue{}%
    \global\let\mermin@tag\df@label
  \fi
  \begingroup
    \let\mermin@found\@empty
    \def\do##1{\def\mermin@tmp{##1}%
      \ifx\mermin@tmp\mermin@form \def\mermin@found{x}\fi}%
    \mermin@queue
    \ifx\mermin@found\@empty
      \toks@\@xp{\mermin@form}%
      \xdef\mermin@queue{%
        \unexpanded\@xp{\mermin@queue}\noexpand\do{\the\toks@}}%
    \fi
  \endgroup
}
%    \end{macrocode}
% \end{macro}
%
% \begin{macro}{\mermin@also@in@display}
% \begin{macro}{\mermin@also@stored}
% \cs{alsolabel} in a display converts its key to stored form.
% If the line has no label yet, the key becomes the primary label and is
% marked as promoted; otherwise it is queued.
%    \begin{macrocode}
\def\mermin@also@in@display#1{%
  \mermin@canonical{#1}%
  \@xp\mermin@also@stored\@xp{\mermin@form}%
}
\def\mermin@also@stored#1{%
  \ifx\df@label\@empty
    \gdef\df@label{#1}%
    \global\let\mermin@promoted\df@label
    \global\let\mermin@tag\df@label
    \gdef\mermin@queue{}%
  \else
    \gdef\mermin@form{#1}%
    \mermin@enqueue
  \fi
}
%    \end{macrocode}
% \end{macro}
% \end{macro}
%
% \begin{macro}{\mermin@real@in@display}
% A real \cs{label} normally passes straight through, and so still reports
% a second \cs{label} on one line as an error.
% If instead the line's primary label was promoted from \cs{alsolabel}, the
% real key takes the slot and the promoted key joins the queue.
% A promotion left over from a discarded measuring pass no longer matches
% |\df@label| and is cleared.
%    \begin{macrocode}
\def\mermin@real@in@display#1{%
  \ifx\mermin@promoted\@empty
    \mermin@call@orig{#1}%
  \else\ifx\mermin@promoted\df@label
    \global\let\mermin@form\df@label
    \ifx\mermin@tag\df@label\else\gdef\mermin@queue{}\fi
    \global\let\df@label\@empty
    \mermin@call@orig{#1}%
    \global\let\mermin@tag\df@label
    \mermin@enqueue
    \gdef\mermin@promoted{}%
  \else
    \gdef\mermin@promoted{}%
    \mermin@call@orig{#1}%
  \fi\fi
}
%    \end{macrocode}
% \end{macro}
%
% \begin{macro}{\mermin@label@in@display}
% A real \cs{label} in a display arrives here, with \textsf{cleveref}'s
% optional argument passed through untouched.
%    \begin{macrocode}
\def\mermin@label@in@display{%
  \@ifnextchar[\mermin@label@in@display@opt\mermin@label@in@display@noarg}
\def\mermin@label@in@display@opt[#1]#2{\mermin@real@in@display{[#1]{#2}}}
\def\mermin@label@in@display@noarg#1{\mermin@real@in@display{{#1}}}
%    \end{macrocode}
% \end{macro}
%
% \begin{macro}{\ltx@label}
% \begin{macro}{\label@in@display}
% \textsf{hyperref} and \textsf{cleveref} install their own |\ltx@label|
% and |\label@in@display| in |\AtBeginDocument| code.
% The package wraps them in the |begindocument/end| hook, which runs after
% all such code, so that load order does not matter.
% The new |\ltx@label| writes the primary label as before and then, if
% the queue belongs to it, every queued key.
% The new |\label@in@display| sends \cs{alsolabel} and \cs{label} their
% separate ways.
%    \begin{macrocode}
\AddToHook{begindocument/end}{%
  \let\mermin@ltx@label\ltx@label
  \def\ltx@label#1{%
    \mermin@ltx@label{#1}%
    \def\mermin@tmp{#1}%
    \ifx\mermin@tmp\mermin@tag
      \begingroup
        \let\do\mermin@ltx@label
        \mermin@queue
      \endgroup
      \gdef\mermin@queue{}\gdef\mermin@tag{}\gdef\mermin@promoted{}%
    \fi}%
  \let\mermin@orig@label@in@display\label@in@display
  \ifdefined\cref@old@label@in@display
    \let\mermin@cref@old\cref@old@label@in@display
  \fi
  \def\label@in@display{%
    \ifmermin@also
      \expandafter\mermin@also@in@display
    \else
      \expandafter\mermin@label@in@display
    \fi}%
}
%</package>
%    \end{macrocode}
% \end{macro}
% \end{macro}
%
% \Finale
\endinput
